% ============================================================
% CHAPTER 3
% ============================================================

\chapter{Computational Expressiveness of SpherePop}
\label{chap:computational-expressiveness}

\section{Motivation}

The SpherePop calculus was originally introduced as a geometric model
of computation based on two primitive operations: \emph{merge} and
\emph{collapse}. Merge implements homotopy colimits (cf.
\cref{chap:merge-hocolim}), whereas collapse abstracts internal detail
while preserving boundary invariants (cf. \cref{chap:collapse-truncation}).
In this chapter we prove that these two operations, together with
geometrically admissible boundary data, suffice to encode universal
computation. In particular, SpherePop is Turing complete.

From a conceptual perspective, this establishes SpherePop as a
fundamental model of geometric computation analogous to the lambda
calculus for functional computation or Turing machines for classical
discrete computation. From the lamphrodynamic point of view, it implies
that computational processes can be realized as lamphrodynamically
admissible evolutions of geometric states.


\section{Encoding of Lambda Calculus}

Let $\Lambda$ denote the usual untyped lambda calculus with application
$(MN)$ and abstraction $(\lambda x.M)$. We construct a geometric encoding
\[
  \Enc \colon \Lambda \longrightarrow \dConfig,
\]
mapping lambda terms to derived configuration objects. Application is
encoded by merge of configurations, and abstraction is encoded by an
entropy--admissible boundary structure which binds a formal variable.

\begin{definition}
Given a lambda term $M$, define $\Enc(M)$ to be the configuration
constructed inductively as follows:
\begin{enumerate}
\item If $M$ is a variable $x$, then $\Enc(x)$ is the basic configuration
representing the boundary label $x$.
\item If $M = (MN)$, then $\Enc(M)$ is the merge
$\Enc(M) \merge \Enc(N)$.
\item If $M = \lambda x.M'$, then $\Enc(M)$ is obtained by collapse of
$\Enc(M')$ along the admissible boundary corresponding to $x$.
\end{enumerate}
\end{definition}

\begin{remark}
This encoding is reminiscent of categorical semantics of the lambda
calculus in cartesian closed categories, except that merge plays the
role of application and collapse plays the role of abstraction. The
boundary structure supplied by lamphrodynamics serves as a binder.
\end{remark}


\section{Beta Reduction as Geometric Evolution}

The key observation is that beta reduction in the lambda calculus
corresponds to a geometric sequence of merge and collapse operations.
Given $(\lambda x.M)N$, the merge operation combines the configurations
$\Enc(M)$ and $\Enc(N)$. A subsequent entropy--admissible collapse
substitutes the boundary condition for $x$ with that of $N$.

\begin{proposition}
Let $M,N$ be lambda terms. Then
\[
  \Enc((\lambda x.M)N) \simeq
  \Collapse(\Enc(M)\merge\Enc(N)),
\]
where the collapse substitutes boundary data associated to $x$.
\end{proposition}

\begin{proof}[Sketch]
Merge produces a configuration in which the boundary associated to $x$
is identified with the boundary associated to $N$. Entropy--admissible
collapse retracts the resulting configuration to one in which $x$ is
substituted by $N$. Details follow the inductive construction of $\Enc$.
\end{proof}

\begin{corollary}
$\Enc$ preserves beta reduction: if $M\to_\beta M'$, then
$\Enc(M)\rightsquigarrow\Enc(M')$ via merge--collapse evolution.
\end{corollary}


\section{Turing Completeness}

We now prove the main theorem of this chapter.

\begin{theorem}[SpherePop Turing Completeness]
\label{thm:spherepop-turing}
The SpherePop calculus is Turing complete. More precisely, the encoding
$\Enc$ is sound and complete with respect to lambda calculus reduction.
Hence any partial recursive function can be represented as a
lamphrodynamically admissible sequence of merge and collapse operations.
\end{theorem}

\begin{proof}[Sketch of proof]
Soundness follows from the previous corollary. For completeness, we
construct a geometric evaluation strategy by induction on the lambda
term, using merge for application and collapse for abstraction. The
fundamental step is to show that every lambda term can be realized as a
finite lamphrodynamic composition of merge and collapse. Since lambda
calculus is Turing complete, so is SpherePop. Full technical details
require a careful treatment of boundary invariants and lamphrodynamic
admissibility, deferred to Appendix~M.
\end{proof}


\section{Complexity Analysis}

Computational complexity in SpherePop depends on the number of merge and
collapse operations as well as the lamphrodynamic cost of collapse.
While complexity in lambda calculus is measured by beta reductions,
complexity in SpherePop is measured by geometric operations.

\begin{definition}
The \emph{SpherePop complexity} $\Comp(X)$ of an encoded configuration
$X=\Enc(M)$ is the minimal number of merge and collapse operations
required to produce a normal form of $X$.
\end{definition}

\begin{proposition}
There exists a constant $C>0$ such that for any lambda term $M$,
\[
  \Comp(\Enc(M)) \le C \cdot \beta(M),
\]
where $\beta(M)$ is the usual beta--reduction complexity of $M$.
\end{proposition}

\begin{proof}[Idea]
Each beta reduction corresponds to at most one merge and one collapse.
The constant $C$ accounts for auxiliary geometric normalizations.
\end{proof}

\begin{remark}
This estimate shows that SpherePop does not incur asymptotic complexity
penalties relative to the lambda calculus, up to constant factors. A
precise characterization of SpherePop complexity classes remains open.
\end{remark}


\section{Comparison with Other Models}

SpherePop differs from standard models of computation in two essential
ways: (i) its primitives are geometric rather than symbolic, and (ii)
computational steps are lamphrodynamically admissible operations. In
particular, SpherePop computation can occur \emph{in situ} in a physical
system governed by lamphrodynamic evolution.

\begin{remark}
Unlike cellular automata or graph rewriting systems, SpherePop is based
on derived geometry and shifted symplectic structures. This makes it a
candidate model for computation in physical theories with geometric
constraints, including RSVP itself.
\end{remark}


\section{Open Problems}

\begin{enumerate}
\item Determine SpherePop complexity classes and their relationship to
standard complexity classes (P, NP, etc.).
\item Develop a geometric type theory for SpherePop analogous to typed
lambda calculus.
\item Investigate lamphrodynamic resource constraints for physical
computation in RSVP.
\end{enumerate}

\medskip
\noindent
This completes Part~I, Chapter~3. The next part develops semantic and
geometric evolutions in L--systems.
\newpage

